% ============================================================
\section{Worked Example: Italian Alps Test Site}
\label{sec:example}
% ============================================================

To illustrate the pipeline output, we apply it to a candidate
installation site in the Italian Alps: a sub-alpine meadow at
approximately $46.4^\circ\,$N, $11.3^\circ\,$E, elevation
\SI{1850}{m\,a.s.l.}
The site is representative of the challenging conditions that
motivate careful pre-installation characterisation: significant
relief, mixed land cover (meadow, sparse conifers, rocky outcrops),
seasonal snow, and clay-rich soils derived from carbonate--flysch
geology.

\subsection{Terrain and Topographic Correction}

Figure~\ref{fig:dem_overview} shows the 30-m Copernicus DEM
mosaicked over the $\pm\SI{500}{m}$ bounding box around the sensor
location.
The terrain is characterised by a north-facing slope
($\sim 8^\circ$ mean slope within $\rzz$) with a prominent ridge to
the south-east rising approximately \SI{120}{m} above the sensor.

\begin{figure}[h]
  \centering
  \includegraphics[width=0.8\linewidth]{dem_hillshade}
  \caption{Hillshaded DEM with the $\rzz$ (inner circle) and
    $3\,\rzz$ (outer circle) footprint boundaries overlaid.
    The sensor location is marked with a cross.}
  \label{fig:dem_overview}
\end{figure}

The 3-D ray-casting algorithm over $18 \times 72$ directions gives
\begin{equation*}
  \kT = 0.97,
\end{equation*}
indicating that the south-eastern ridge reduces the effective soil
volume by approximately 3\% relative to a flat reference.
Without this correction, the Desilets $N_0$ would be overestimated
by the same margin.

\subsection{Muon Field-of-View Correction}

The horizon profile $\psi(\phi)$ shows a maximum elevation angle of
$\SI{8.5}{\degree}$ towards azimuth $\SI{145}{\degree}$
(south-south-east) and is below $\SI{3}{\degree}$ for all other
azimuths.
The analytical $\kM$ integral (Eq.~\eqref{eq:kappa_muon}) gives
\begin{equation*}
  \kM = 0.992,
\end{equation*}
a $0.8\%$ reduction in the effective muon flux.
This is small but non-negligible at a site with high-precision
calibration requirements.

\subsection{Land-Use / Land-Cover Correction}

Figure~\ref{fig:lulc_map} shows the ESA WorldCover 2021 map
reprojected to the DEM grid.
Within the $\rzz$ circle the land cover is dominated by grassland
($f_H = 0.95$, \SI{64}{\%} of weighted area), tree cover
($f_H = 1.50$, \SI{18}{\%}), and bare rock
($f_H = 0.03$, \SI{11}{\%}).
A small seasonal stream mapped in OSM ($f_H = 4.0$) covers
$\SI{0.4}{\%}$ of the $W(r)$-weighted area.

\begin{figure}[h]
  \centering
  \includegraphics[width=0.8\linewidth]{lulc_map}
  \caption{ESA WorldCover 2021 classification within the
    $2\,\rzz$ footprint circle, with OSM water features overlaid.
    The $\rzz$ boundary is the solid white circle.}
  \label{fig:lulc_map}
\end{figure}

The weighted $\kL$ (Eq.~\eqref{eq:kappa_lulc}) evaluates to
\begin{equation*}
  \kL = 1.08,
\end{equation*}
reflecting the net positive hydrogen contribution from tree canopy
water relative to the bare-soil baseline.
Without this correction, the calibrated $N_0$ would be
overestimated by $\sim 8\%$, propagating directly into a
soil-moisture bias of $\sim \SI{0.02}{\cubic\metre\per\cubic\metre}$.

\subsection{Total Correction and Reference Count}

The total correction factor is
\begin{equation*}
  \kTot = \kT \cdot \kM \cdot \kL
        = 0.97 \times 0.992 \times 1.08
        = 1.039.
\end{equation*}
With a detector-specific sea-level reference count
$N_\mathrm{sl} = \SI{2800}{\text{cph}}$ and a site rigidity
cutoff $R_c = \SI{4.2}{\giga\volt}$ (retrieved from \texttt{crnpy}):
\begin{align*}
  f_{R_c} &= \exp[\alpha_H (R_c^\mathrm{site} - R_c^\mathrm{ref})]
           = \exp[-0.075 \times (4.2 - 4.49)] = 1.022, \\
  N_\mathrm{neut,site} &= 2800 \times 1.022
    \times \exp[0.0077 \times (1013.25 - 794)] \times 1.039
    = \SI{3620}{\text{cph}}.
\end{align*}
With $\lw = 0.055$ g/g (from SoilGrids clay fraction
$f_\mathrm{clay} = 0.23$, Eq.~\eqref{eq:lw_clay}):
\begin{equation*}
  N_0 = \frac{3620}{0.0808 + 0.372\,/\,(0.055 + 0.115)}
      = \frac{3620}{0.0808 + 2.188}
      \approx \SI{1601}{\text{cph}}.
\end{equation*}
This value enters the Desilets inversion
(Eq.~\eqref{eq:desilets_inv}) for all subsequent soil-moisture
retrievals.

\subsection{Soil Properties and Supplementary Hydrogen Inventory}

Table~\ref{tab:soilprops} summarises the SoilGrids-derived soil
properties for the 0–30 cm layer.

\begin{table}[h]
\centering
\caption{Soil properties at the Italian Alps test site (0–30 cm
depth-weighted mean, SoilGrids 2.0).}
\label{tab:soilprops}
\begin{tabular}{lll}
\toprule
Property & Value & Unit \\
\midrule
Bulk density $\rho_b$          & 1.18  & \si{g\,cm^{-3}} \\
Clay fraction $f_\mathrm{clay}$ & 0.23  & g/g \\
SOC content                    & 28    & \si{g\,kg^{-1}} \\
Lattice water $\lw$            & 0.055 & g/g \\
$\theta_\mathrm{SOC}$          & 0.028 & \si{m^3\,m^{-3}} \\
$z_{86}$ (at $\thetav = 0.20$) & 26    & cm \\
\bottomrule
\end{tabular}
\end{table}

\subsection{Atmospheric and Vegetation Seasonal Climatology}

Figure~\ref{fig:seasonal} shows the monthly climatological
$\rho_{WV}$, $f_{WV}$, SWE and LAI averaged over 2019–2023.

\begin{figure}[h]
  \centering
  \includegraphics[width=\linewidth]{seasonal_climatology}
  \caption{Monthly climatology at the Italian Alps test site.
    \textit{Top left}: atmospheric water-vapour density
    $\rho_{WV}$ [\si{g\,m^{-3}}].
    \textit{Top right}: neutron correction factor $f_{WV}$.
    \textit{Bottom left}: snow-water equivalent SWE [mm].
    \textit{Bottom right}: MODIS LAI.}
  \label{fig:seasonal}
\end{figure}

Key observations:
\begin{itemize}[nosep]
  \item $\rho_{WV}$ ranges from $\sim\SI{3}{\gram\per\cubic\metre}$
        in January to $\sim\SI{12}{\gram\per\cubic\metre}$ in July,
        implying an operational $f_{WV}$ variation of 5–6\% between
        winter and summer.
  \item SWE exceeds the \SI{150}{mm} opacity threshold in January–March,
        making those months unsuitable for soil-moisture calibration.
  \item Peak LAI in July ($\approx 3.2$ m$^2$/m$^2$) gives
        $\mathrm{AGBH} \approx \SI{480}{\milli\metre}$,
        equivalent to $\sim\SI{0.048}{\cubic\metre\per\cubic\metre}$
        additional hydrogen — an important bias term in summer.
\end{itemize}

\subsection{Optimal Sampling Plan}

Figure~\ref{fig:sampling} shows the recommended 26-point sampling
plan superimposed on the DEM and the $\rzz(\thetav)$ dynamic-range
curve.

\begin{figure}[h]
  \centering
  \includegraphics[width=\linewidth]{sampling_plan}
  \caption{\textit{Left}: recommended soil sampling locations
    (filled circles, coloured by ring) overlaid on the hillshaded DEM.
    Ring radii are $\{0.25, 0.50, 0.75, 1.00\}\times\rzz$.
    \textit{Right}: footprint radius $\rzz$ as a function of
    volumetric soil moisture $\thetav$ at site pressure.
    Dashed line marks the design moisture $\thetav = 0.20$.}
  \label{fig:sampling}
\end{figure}

At the design moisture of
$\thetav = \SI{0.20}{\cubic\metre\per\cubic\metre}$,
$\rzz = \SI{113}{\metre}$.
The four sampling rings are therefore located at
$r = $ \SI{28}, \SI{57}, \SI{85}, \SI{113} m.
The inner ring (4 points) carries $\Phi_1 / \sum \Phi = 18\%$ of
the total footprint weight; the outer ring (8 points) carries 32\%.
The dry-soil extreme ($\thetav \to 0$) gives
$\rzz \to \SI{188}{\metre}$; at saturation
($\thetav = 0.45$) $\rzz$ shrinks to $\SI{106}{\metre}$,
a 44\% dynamic range.
